\section{Experimental setup}
\label{sec:setup}

\subsection{Why a synthetic primary benchmark}

Section~\ref{sec:external} reports the public dataset survey in full; the conclusion
directly motivates this section. No public cybersecurity benchmark is designed for
continuous Zero Trust authorization, and none provides principal identity alongside
client TLS fingerprints. A corpus populating the 5-node chain at scale with labeled
lateral movement and credential reuse does not currently exist in the literature. A synthetic streaming
environment is therefore constructed, and its statistical properties are systematically
validated, to eliminate trivial separability shortcuts.

\subsection{The generator}

The generator simulates the chain \textit{ip} $\rightarrow$ \textit{config (JA3)}
$\rightarrow$ \textit{device} $\rightarrow$ \textit{user} $\rightarrow$ \textit{resource}
over a population of users, devices, sources, client configurations, and procedurally
generated resources, with circadian arrival rates, roaming and NAT, shared devices,
cookie wipes, periodic service account batch traffic, and benign human error.
It emits binary anomaly labels, per-class type vectors, and a benign-context bitmask.

Benign principals perform legitimate non-habitual accesses with non-zero
probability. Without exploratory benign traffic, lateral movement would trivially coincide
with destination novelty, allowing models to achieve high performance without modeling temporal
dynamics. By introducing benign exploration, an individual lateral event cannot be separated
by destination novelty alone, requiring models to discriminate based on relational history
and temporal sequences.

\subsection{The leakage audit}

\begin{table}[ht!]
\centering
\renewcommand{\arraystretch}{1.5}
\begin{tabular}{lcc}
\toprule
Class & Before & After \\
\midrule
Policy      & \FloorPolicyPre & \FloorPolicyPost~\textit{(risk; semantic, declared)} \\
Contextual  & \FloorCtxPre    & \FloorCtxPost~\textit{(sensors; by design)} \\
\textit{Lateral}    & \textit{\FloorLatPre} & \textit{\FloorLatPost} \\
Credential reuse & \FloorTheftPre & \FloorTheftPost \\
\bottomrule
\end{tabular}
\caption{Best single-feature AUC per class, before and after de-leakage. A high value
indicates that a single input column separates the class on its own, trivializing the task
and obscuring model evaluation. Post-de-leakage values from
\texttt{tasks/runs/leakage\_audit\_floor.log}; pre-de-leakage values are from the earlier
audit on the enriched (pre-de-leak) generator.}
\label{tab:floor}
\end{table}

In synthetic security benchmarks, unconstrained parameter spaces frequently introduce
subtle artifact leakage. For instance, generating resources under uniform attack sampling
while benign traffic follows a power law causes resource identifiers to become an artificial
shortcut (\FloorLatPre{} single-feature AUC). Similarly, response-side features (such as HTTP
status codes or payload volume) act as direct label proxies unavailable at decision time.

To eliminate these artifacts, strict synthetic invariance is enforced: resource popularity
is decoupled from indexing via seed-dependent permutations, attack destinations are drawn
from the identical marginal distribution as benign traffic, deterministic categorical signatures
are replaced with continuous variations, and response-side fields are strictly excluded.
As shown in Table~\ref{tab:floor}, the single-feature baseline AUC on lateral movement drops
from \FloorLatPre{} to \FloorLatPost{}, establishing that no individual feature trivially
separates the attack classes.

These constraints are formalized into an automated leakage audit suite executed across
random seeds. The audit verifies: (1) an upper bound on single-feature AUC; (2) absence of exact
constant fingerprints; (3) statistical equivalence between benign and anomalous destination
marginal distributions (evaluated by effect size rather than asymptotic $p$-values, avoiding
spurious rejections at large sample sizes); and (4) valid routing semantics for all emitted
requests. Lateral movement and credential reuse are subject to no single-feature
leakage exemptions.

\subsection{Protocol}

Unless stated otherwise, experiments follow the standard configuration (\PanelAsettings).
Training is conducted on benign events only. The test stream is replayed chronologically,
event by event, under the external commit gate described in Section~\ref{sec:problem}.
Dispersion is reported as the sample standard deviation over the 3 seeds, 1 run
per seed, with $\text{ddof}=1$; paired comparisons use per-seed differences evaluated with a
Wilcoxon signed-rank test and bootstrap confidence intervals. With 3 seeds the signed-rank
test has a minimum attainable $p$ of $0.25$, so the sign pattern of the
per-seed differences is additionally reported, and significance is not claimed from the test alone.

\subsection{Baselines}

5 baselines are evaluated alongside the model: Isolation Forest, One-Class SVM,
XGBoost, a static graph neural network (identical graph structure without temporal memory),
and a standard 2-node bipartite TGN (\textit{user} $\rightarrow$ \textit{resource}).
All baselines receive the same per-event tabular signals as the TGN: the edge message
features, the static attributes of both endpoints, causal benign-gated interaction
counters, and the identical precursor prior. The baselines consume these as a flat
per-event vector for the device actor, so the comparison isolates the specific
contribution of temporal graph modeling.
The static GNN isolates the role of recurrent temporal state, while the 2-node TGN isolates
the effect of the 5-node decomposition against conventional bipartite link prediction.
XGBoost is fully supervised (trained on ground-truth labels) and is included solely as an
empirical upper-bound reference. For the one-class Isolation Forest, hyperparameters are
selected on the validation slice by validation AUC (a standard compromise of the one-class
setting, in which no labels enter the training fit itself).
